% Part: normal-modal-logic
% Chapter: axioms-systems
% Section: logics-proofs

\documentclass[../../../include/open-logic-section]{subfiles}

\begin{document}

\olfileid{nml}{prf}{prf}

\olsection{\usetoken{P}{derivation} आणि मोडल प्रणाली}

सामान्य मोडल तर्कशास्त्रांसाठी !!a{derivation}
म्हणजे काय ते आधी ठरवू. साधारणपणे,
!!a{derivation} म्हणजे !!{formula}s ची अशी मालिका
की तिच्यातील प्रत्येक !!{element} काही
\emph{स्वयंसिद्धकांपैकी} एक (किंवा त्याचा आदेशन-दृष्टांत)
असतो, किंवा आधीच्या !!{element}s पासून मोजक्या
अनुमान नियमांपैकी एखाद्याने मिळतो. सामान्य मोडल
तर्कशास्त्रात सर्व सर्वतः-सत्य सूत्रांचे दृष्टांत
\iftag{prvDiamond}{, \Ax{K}, आणि~\Dual{}}{ आणि~\Ax{K}}
स्वयंसिद्धके मानली जातात. त्यातून लघुतम सामान्य
मोडल तर्कशास्त्र, म्हणजे मोडल प्रणाली~$\Log{K}$,
मिळते. इतर प्रणाली मिळवण्यासाठी आणखी स्वयंसिद्धके
जोडता येतात. नियम मात्र नेहमी विध्यात्मक अनुमान~\MP{}
आणि अनिवार्यीकरण~\Nec हेच असतात.

\begin{defn}
  मोडल प्रणाली $\Log{K} !A_1 \dots !A_n$ आणि
  !!a{formula}~$!B$ दिले असता, $!C_1$, \dots,
  $!C_k$ अशी !!{formula}s असतील, जिथे $!C_k = !B$
  आणि प्रत्येक $!C_i$ हा सर्वतः-सत्य दृष्टांत,
  किंवा $\Ax{K}$,\iftag{prvDiamond}{ $\Dual$,}{}
  $!A_1$, \dots, $!A_n$ यांपैकी एखाद्याचा दृष्टांत,
  किंवा आधीच्या !!{formula}s पासून \MP{} अथवा~\Nec
  नियमाने मिळणारे सूत्र असेल, तर आणि तेव्हाच
  $!B$ हे $\Log{K} !A_1 \dots
  !A_n$ मध्ये \emph{!!{derivable}} आहे असे म्हणतो;
  संकेतात $\Log{K} !A_1 \dots !A_n \Proves !B$.
\end{defn}

पुढील प्रतिज्ञेनुसार $!B$ ची $\Sigma$-!!{derivation}
दाखवून $!B \in \Sigma$ हे सिद्ध करता येते.

\begin{prop}
  $\Log{K} !A_1 \dots !A_n = \Setabs{!B}{\Log{K} !A_1 \dots !A_n \Proves !B}$.
\end{prop}

\begin{proof}
  !!{derivation}s च्या लांबीवर विगमन करून
  $\Setabs{!B}{\Log{K} !A_1 \dots !A_n \Proves !B} \subseteq
  \Log{K} !A_1 \dots !A_n$ हे दाखवू.

  $!B$ च्या !!{derivation} ची लांबी~$1$ असेल,
  तर त्यात एकच !!{formula} असते. ते !!{formula} आधीच्या
  सूत्रांपासून \MP{} किंवा \Nec ने मिळालेले
  असू शकत नाही. म्हणून ते सर्वतः-सत्य दृष्टांत,
  \Ax{K} चा,\iftag{prvDiamond}{ $\Dual$ चा,}{}
  किंवा $!A_1$, \dots,~$!A_n$ यांपैकी एखाद्याचा
  दृष्टांत असला पाहिजे. $\Log{K}!A_1\dots!A_n$
  हे सर्व सामावते; म्हणून $!B \in \Log{K}!A_1
  \dots !A_n$.

  $!B$ च्या !!{derivation} ची लांबी~$> 1$
  असेल, तर आधीच्या प्रकरणांखेरीज !!{derivation}
  मधील अंतिम ओळ नसलेल्या !!{formula}s पासून
  \MP{} किंवा \Nec{} नेही
  $!B$ मिळू शकते. $!B$ हे $!C$ आणि
  $!C \lif !B$ यांपासून (\MP ने) मिळत असेल,
  तर विगमन गृहीतकावरून $!C$ आणि
  $!C \lif !B \in
  \Log{K}!A_1 \dots !A_n$. प्रत्येक मोडल
  तर्कशास्त्र विध्यात्मक अनुमानाखाली बंद असते;
  म्हणून $!B \in \Log{K}!A_1 \dots
  !A_n$. $!B \equiv \Box !C$ हे $!C$
  पासून \Nec ने मिळत असेल, तर विगमन
  गृहीतकावरून $!C
  \in \Log{K}!A_1 \dots !A_n$. प्रत्येक
  सामान्य मोडल तर्कशास्त्र $\Nec$ खाली बंद असते;
  म्हणून $!B \in
  \Log{K}!A_1\dots!A_n$.

  उलट समावेशासाठी $\Sigma = \Setabs{!B}{\Log{K} !A_1 \dots !A_n \Proves !B }$
  हे $!A_1$, \dots, $!A_n$ चे सर्व दृष्टांत
  सामावणारे सामान्य मोडल तर्कशास्त्र आहे, हे
  दाखवू. मग $\Log{K} !A_1 \dots !A_n$
  हे व्याख्येने असे लघुतम तर्कशास्त्र आहे,
  या तथ्याचा उपयोग करू.
  \begin{enumerate}
    \item प्रत्येक सर्वतः-सत्य सूत्र~$!B$ हाही
      सर्वतः-सत्य दृष्टांत आहे. म्हणून
      $\Log{K}!A_1\dots!A_n \Proves !B$;
      त्यामुळे $\Sigma$ सर्व सर्वतः-सत्य सूत्रे सामावते.
    \item $\Log{K}!A_1\dots!A_n \Proves !C$
      आणि $\Log{K}!A_1\dots!A_n \Proves !C \lif !B$
      असल्यास $\Log{K}!A_1\dots!A_n \Proves !B$:
      $!C$ ची !!{derivation} आणि $!C \lif !B$
      ची !!{derivation} जोडून अंतिम ओळ~$!B$
      घाला. या ओळीला \MP{} चे समर्थन मिळते.
      म्हणून $\Sigma$ विध्यात्मक अनुमानाखाली बंद आहे.
    \item $!B$ ची !!a{derivation} असेल, तर
      $!B$ च्या प्रत्येक आदेशन-दृष्टांताचीही
      !!a{derivation} असते: संपूर्ण !!{derivation}
      मधील प्रत्येक !!{formula} वर तेच आदेशन लावा.
      (सराव: !!{derivation}s च्या लांबीवर विगमन
      करून परिणामी मालिकाही योग्य निष्पत्ती
      आहे हे सिद्ध करा.) म्हणून $\Sigma$ एकरूप
      आदेशनाखाली बंद आहे. ($\Sigma$ मोडल तर्कशास्त्राच्या
      सर्व अटी पूर्ण करते, हे आता सिद्ध झाले.)
    \item $\Log{K}!A_1\dots!A_n \Proves \Ax{K}$;
      म्हणून $\Ax{K} \in \Sigma$.
    \tagitem{prvDiamond}{$\Log{K}!A_1\dots!A_n \Proves \Dual$
      असल्यामुळे $\Dual \in \Sigma$.}{}
    \item $\Log{K}!A_1\dots!A_n \Proves !C$
      असेल, तर अतिरिक्त ओळ~$\Box !C$ ला~\Nec
      चे समर्थन मिळते. त्यामुळे $\Sigma$ हा
      \Nec खाली बंद आहे; म्हणजे $\Sigma$ सामान्य आहे.
  \end{enumerate}
\end{proof}

\end{document}
